%==============================================================================
% Section 3: Hamiltonian admissibility by topology
%==============================================================================
\section{Hamiltonian admissibility by topology}
\label{sec:hamiltonian}

\subsection{Admissibility criterion}
\label{sec:hamiltonian:criterion}

In this paper, admissibility in the reduced sector is judged by the
following three criteria.

\begin{enumerate}
  \item \emph{EH constraint closure}: the Hamiltonian constraint of the pure
        $\EH$ part closes under the Lorentzian convention
        (\code{EH_PASS}).
  \item \emph{Reduced action skeleton}: the EC+NY reduced Lagrangian of each
        torsion mode has a well-defined velocity structure
        (\code{LRED_PASS}).
  \item \emph{Torsion auxiliary closure}: the algebraic equations for the
        auxiliary variables $\eta$ and $V$ admit real solutions
        (\code{AUX_PASS} or \code{AUX_CONDITIONAL_REAL_BRANCH}).
\end{enumerate}

This notion of admissibility is not a claim of orbit stability, unitarity,
or full constraint closure.

\subsection{Hamiltonian constraint}
\label{sec:hamiltonian:constraint}

Lapse variation $\delta S/\delta N = 0$ gives the Hamiltonian constraint
\begin{equation}
  \Ham
  \equiv A_{\topo}\, q\!\left(F_{\topo}+\frac{\dot q^{2}}{N^{2}}\right) = 0,
  \qquad
  \text{that is,}\quad F_{\topo}+\frac{\dot q^{2}}{N^{2}} = 0.
\end{equation}
On the $\EH$ branch ($\eta=V=0$) one has $F_{\topo}\big|_{\EH}=\chi$ (with
the $\chi$ of Table~\ref{tab:topology_chi} substituted for each topology),
which is the content of \code{EH_PASS}:

\begin{table}[H]
\centering
\begin{tabular}{lcc}
\toprule
Topology & $F_{\EH}$ (on-shell) & Constraint residual \\
\midrule
$\Sthree$   & $4$      & $4N^{2}+\dot r^{2}=0$ \\
$\Tthree$   & $0$      & $\dot r^{2}=0$ \\
$\Nilthree$ & $-1/12$  & $N^{2}-12\dot R^{2}=0$ \\
$\Solthree$ & $-1/3$   & $3N^{2}-9\dot q^{2}=0$\ (local/iso) \\
\bottomrule
\end{tabular}
\caption{$\EH$ residuals for each topology.}
\label{tab:eh_residual}
\end{table}

The torsion auxiliary equations are the algebraic equations
$\delta S/\delta\eta=0$ and $\delta S/\delta V=0$.  On $\AX$ only $\eta$ is
non-trivial and on $\VT$ only $V$, each yielding \code{AUX_PASS}.  On $\MX$
the auxiliary Hessian determinant is
\begin{equation}
  \det\!\begin{pmatrix}
    \partial_{\eta}(\partial_{\eta} F) & \partial_{V}(\partial_{\eta} F)\\
    \partial_{\eta}(\partial_{V} F)    & \partial_{V}(\partial_{V} F)
  \end{pmatrix}
  = -\frac{4q^{2}(\kap^{4}\thetaNY^{2}+1)}{9} < 0.
\end{equation}
For $q>0$ in the reduced-sector domain and real $\kap$ and $\thetaNY$, one has
$\kap^{4}\thetaNY^{2}+1>0$, so the auxiliary Hessian is non-degenerate.
We retain this mixed auxiliary channel as a branch with real-branch
selection and origin tagging, and record it as
\code{AUX_CONDITIONAL_REAL_BRANCH}, that is, as the
\code{H_CONDITIONAL} status of $\MX$.

The negative determinant does not imply a stable minimum.  The notion of
\emph{conditional admissibility} used here is a weak criterion requiring
only algebraic closure and real-branch selection at the non-degenerate
mixed auxiliary saddle, and is independent of the Hessian signature
(dynamical stability).

\subsection{Mode status}
\label{sec:hamiltonian:mode}

Collecting the above criteria gives Table~\ref{tab:ham_status}.  The result
is the same across the four topologies.

\begin{table}[H]
\centering
\begin{tabular}{lll}
\toprule
Mode & Hamiltonian status & Interpretation \\
\midrule
$\EH$ & \code{H_PASS}         & pure $\EH$ reference branch \\
$\AX$ & \code{H_PASS}         & active-torsion admissible \\
$\VT$ & \code{H_PASS}         & active-torsion admissible \\
$\MX$ & \code{H_CONDITIONAL}  & conditional admissibility via real auxiliary branch \\
\bottomrule
\end{tabular}
\caption{Hamiltonian status by topology and mode.}
\label{tab:ham_status}
\end{table}

The \code{H_CONDITIONAL} label for $\MX$ is not an obstruction or a failure.
Assuming the existence of a real auxiliary branch, the off-shell behaviour
of the subsequent P-channel diagnostics and the $\QNY$ boundary channel is
recorded as a diagnostic, which is the content of conditional admissibility.

\paragraph{Remark.}
\code{H_PASS} and \code{L_ADMISSIBLE} do not imply dynamical stability.  The
$\Solthree$ table entries refer to the local/isotropic reduction and do not
constitute a full resolution of the compact quotient.
